% Escribe aquí tu parte. NO toques el main.
% Usa \section{}, \subsection{}, etc.
% Para imágenes usa: \includegraphics[width=\linewidth]{figures/TU_CARPETA/imagen.png}

\chapter{Modelo Clave-Valor (Fer)}
\label{ch:fer}

El modelo \textbf{clave-valor} es una de las aproximaciones NoSQL más simples y eficientes: cada dato lo almacena como un par (\textit{key}, \textit{value}). A diferencia del modelo relacional, no existe un esquema rígido con tablas y columnas, sino un espacio de claves donde se consultan valores mediante accesos directos. Este enfoque resulta especialmente útil en sistemas con alta concurrencia, datos temporales y operaciones muy frecuentes, como los escenarios típicos de una red social.

Para este trabajo se ha seleccionado \textbf{Redis} como SGBD clave-valor de referencia, ya que combina velocidad (principalmente en memoria), facilidad de despliegue y un conjunto de estructuras de datos que permiten modelar relaciones y estados de forma sencilla.

\section{SGBD recomendado: Redis}

Redis (REmote DIctionary Server) es un sistema NoSQL clave-valor orientado a \textbf{rendimiento} y \textbf{baja latencia}, ampliamente utilizado como caché, gestor de sesiones, colas, contadores y almacenamiento temporal. Su valor añadido frente a un almacén clave-valor básico es que no solo almacena cadenas (\texttt{string}), sino también estructuras de datos de alto nivel como \texttt{hashes}, \texttt{sets}, \texttt{lists} o \texttt{sorted sets}, muy útiles para sistemas de interacción social.

\section{Instalación y Configuración del Entorno}

Redis es sencillo de desplegar en Linux y permite comenzar a trabajar en pocos minutos.

\subsection{Instalación en Ubuntu/Debian}
\begin{minted}{bash}
sudo apt update
sudo apt install redis-server
sudo systemctl enable --now redis-server
redis-cli ping
\end{minted}

Si todo es correcto, el comando \texttt{PING} responde con \texttt{PONG}.

\subsection{Instalación mediante Docker (alternativa)}
\begin{minted}{bash}
docker run -d --name redis-ekis -p 6379:6379 redis:latest
redis-cli -h 127.0.0.1 -p 6379 ping
\end{minted}

\section{Diseño Clave-Valor del Subsistema de Usuarios}

En una red social como EKIS, Redis se adapta especialmente bien a la gestión de:
\begin{itemize}
    \item \textbf{Sesiones y tokens} con caducidad (TTL).
    \item \textbf{Relaciones rápidas} (amistades, bloqueos) mediante conjuntos.
    \item \textbf{Contadores} (likes, seguidores) mediante incrementos atómicos.
    \item \textbf{Caché de perfiles} para acelerar lecturas frecuentes.
\end{itemize}

\subsection{Estrategia de claves}
Se define unas reglas simples para garantizar consistencia en el diseño:
\begin{itemize}
    \item \texttt{user:<id>} $\rightarrow$ datos del usuario en \texttt{HASH}
    \item \texttt{friends:<id>} $\rightarrow$ conjunto de amistades en \texttt{SET}
    \item \texttt{blocked:<id>} $\rightarrow$ conjunto de bloqueados en \texttt{SET}
    \item \texttt{session:<token>} $\rightarrow$ sesión activa en \texttt{STRING} con TTL
\end{itemize}

\section{Equivalencias DDL y DML en Redis}

\subsection{DDL (definición de estructura)}
Redis no requiere DDL tradicional (no existen \texttt{CREATE TABLE}). La estructura se define por el tipo de dato asociado a cada clave. En entornos reales, la consistencia del diseño se garantiza por convención de claves y por el código de la aplicación.

\subsection{DML (operaciones típicas)}
Redis se maneja con comandos directos:
\begin{itemize}
    \item \texttt{SET / GET} para valores simples (sesiones, flags).
    \item \texttt{HSET / HGETALL} para perfiles y datos agrupados en \texttt{HASH}.
    \item \texttt{SADD / SISMEMBER / SMEMBERS} para relaciones en \texttt{SET}.
    \item \texttt{INCR} para contadores atómicos.
    \item \texttt{EXPIRE / TTL} para caducidad (sesiones).
\end{itemize}

\section{Implementación y Pruebas}

A continuación se muestran operaciones reales que representan el subsistema de usuarios en EKIS usando \texttt{redis-cli}.

\subsection{Inserción de un usuario (HASH)}
\begin{minted}{bash}
HSET user:10 nombre "Fernando" email "fer@email.com" rol "user"
HSET user:20 nombre "Ismael"   email "isma@ugr.es"   rol "user"
HGETALL user:10
\end{minted}

\subsection{Creación de relación de amistad (SET)}
Para modelar amistad se usan conjuntos. Una amistad bidireccional se representa añadiendo cada usuario al conjunto del otro:

\begin{minted}{bash}
SADD friends:10 20
SADD friends:20 10
SMEMBERS friends:10
\end{minted}

\subsection{Control de coherencia: evitar auto-relación}
Redis no tiene \textit{triggers}. La regla “no puedo ser amigo de mí mismo” se controla desde la lógica de aplicación (o mediante scripts Lua). Por ejemplo, la aplicación tiene que impedir:
\begin{minted}{bash}
SADD friends:10 10
\end{minted}

\subsection{Bloqueo de usuario (SET) con consistencia}
Cuando un usuario bloquea a otro, es coherente eliminar la amistad previa y registrar el bloqueo. Redis nos permite agrupar operaciones con \texttt{MULTI/EXEC}:

\begin{minted}{bash}
MULTI
SREM friends:10 20
SREM friends:20 10
SADD blocked:10 20
EXEC
\end{minted}

\subsection{Gestión de sesión con caducidad (TTL)}
Las sesiones son un caso ideal para Redis porque permite caducidad automática:

\begin{minted}{bash}
SET session:token_abc "10"
EXPIRE session:token_abc 3600
TTL session:token_abc
GET session:token_abc
\end{minted}

\subsection{Ejemplo práctico con Python (redis-py)}
Redis se integra fácilmente con Python para automatizar pruebas y lógica de negocio:

\begin{minted}{python}
import redis

r = redis.Redis(host="127.0.0.1", port=6379, decode_responses=True)

# Crear usuario (HASH)
r.hset("user:10", mapping={"nombre": "Fernando", "email": "fer@email.com", "rol": "user"})

# Añadir amistad (SET)
r.sadd("friends:10", "20")
r.sadd("friends:20", "10")

# Crear sesión con TTL
r.set("session:token_abc", "10")
r.expire("session:token_abc", 3600)

print("Usuario 10:", r.hgetall("user:10"))
print("Amigos de 10:", r.smembers("friends:10"))
print("TTL sesión:", r.ttl("session:token_abc"))
\end{minted}

\section{Adecuación al Sistema EKIS}

El modelo clave-valor con Redis es especialmente adecuado en EKIS como soporte de alto rendimiento para el subsistema de usuarios:

\begin{itemize}
    \item \textbf{Baja latencia y alta concurrencia:} ideal para operaciones muy frecuentes como consultas de estado, contadores y relaciones.
    \item \textbf{Sesiones y expiración automática:} permite gestionar tokens de autenticación y sesiones sin necesidad de procesos de limpieza.
    \item \textbf{Relaciones rápidas con SET:} amistades, bloqueos y listas de seguimiento se resuelven con operaciones eficientes.
    \item \textbf{Soporte natural para caché:} reduce carga en el SGBD relacional principal cuando se consultan perfiles repetidamente.
\end{itemize}
sin embargo, Redis no sustituye completamente al modelo relacional cuando se requieren consultas complejas, \texttt{JOIN}, o garantías transaccionales fuertes a nivel de múltiples entidades.
